\section{Técnicas e instrumentos de recolección de datos}
\label{sec:tecnicas}

La técnica de levantamiento corresponde a la \textbf{simulación computacional
numérica y generación sintética de series temporales}, implementada en Python~3
con NumPy y SciPy. Se emplean cuatro instrumentos operativos, cuyo propósito,
tipo de información producida y estado de avance se detallan a continuación. El
estado declarado no es una afirmación de intención: cada instrumento corresponde
a un módulo de código ejecutable incluido en el anexo digital.

\subsection{Instrumento 1: Generador de referencia física (\emph{truth model})}

\textbf{Propósito.} Producir la trayectoria dinámica real del sistema
($\mathbf{x}_{k}$) libre de perturbaciones de medición, que opera como verdad de
referencia contra la cual se contrasta tanto la telemetría degradada como la
estimación del módulo.

\textbf{Información que permite obtener.} Series temporales del volumen real y
de la tasa real de llenado, en \si{\metre\cubed} y
\si{\metre\cubed\per\second} respectivamente, para los $N$ pasos del horizonte.

\textbf{Forma de aplicación.} Implementa las ecuaciones de espacio de estados de
la sección~\ref{sec:unidad}, propagando el estado mediante las matrices
$\mathbf{A}$ y $\mathbf{B}$ e inyectando la perturbación de proceso
$\mathbf{w}_{k}$ muestreada por factorización de Cholesky de $\mathbf{Q}$.

\textbf{Relación con los objetivos.} Materializa OE2 y provee el denominador de
todas las métricas de error de OE4.

\textbf{Estado.} Implementado, calibrado y verificado.

\subsection{Instrumento 2: Generador de telemetría SCADA degradada}

\textbf{Propósito.} Simular la degradación de la señal del sensor físico bajo
condiciones extremas de vibración y polvo abrasivo, produciendo la lectura
$y_{k}$ que un operador vería en consola.

\textbf{Información que permite obtener.} La serie temporal de mediciones
ruidosas de volumen, que constituye simultáneamente la entrada del filtro y la
\textbf{línea base} contra la cual se mide su aporte.

\textbf{Forma de aplicación.} Aplica la ecuación de medición inyectando ruido
gaussiano $v_{k} \sim \mathcal{N}(0, R)$ sobre la componente observada del
estado real.

\textbf{Relación con los objetivos.} Materializa OE1, al traducir la
caracterización estadística del entorno en una perturbación concreta y
parametrizada.

\textbf{Estado.} Completamente operativo.

\subsection{Instrumento 3: Módulo de software analítico (Filtro de Kalman)}

\textbf{Propósito.} Implementar el ciclo recursivo predictor-corrector que
actualiza la estimación $\hat{\mathbf{x}}_{k|k}$ y minimiza la matriz de
covarianza del error $\mathbf{P}_{k|k}$.

\textbf{Información que permite obtener.} Series temporales del volumen y la
tasa estimados, junto con el historial completo de la ganancia
$\mathbf{K}_{k}$, la covarianza $\mathbf{P}_{k|k}$ y la innovación de cada
ciclo, lo que habilita el análisis de convergencia asintótica.

\textbf{Forma de aplicación.} El núcleo se implementa como clase genérica
desacoplada del proceso físico: recibe las matrices del sistema y opera sin
conocer el dominio de aplicación. La actualización de covarianza emplea la forma
de Joseph,
\[
\mathbf{P}_{k|k} = (\mathbf{I}-\mathbf{K}_{k}\mathbf{H})\,\mathbf{P}_{k|k-1}\,
(\mathbf{I}-\mathbf{K}_{k}\mathbf{H})^{\mathsf{T}}
+ \mathbf{K}_{k}\mathbf{R}\mathbf{K}_{k}^{\mathsf{T}},
\]
algebraicamente equivalente a la forma simplificada
$(\mathbf{I}-\mathbf{K}\mathbf{H})\mathbf{P}$ pero numéricamente estable:
preserva la simetría y la definición positiva de $\mathbf{P}$ frente a errores
de redondeo en aritmética de punto flotante, condición necesaria para
horizontes largos.

\textbf{Relación con los objetivos.} Materializa OE3.

\textbf{Estado.} Implementado y verificado.

\subsection{Instrumento 4: Registro estructurado de trazabilidad}

\textbf{Propósito.} Serializar cada variable temporal con su sello de paso $k$,
garantizando auditoría completa del experimento por terceros.

\textbf{Información que permite obtener.} Archivos CSV con dieciséis columnas
por cada uno de los \num{1000} pasos de la corrida de referencia —entrada de
control, estado real, lectura degradada, estimación, errores de ambas fuentes,
componentes de la ganancia, diagonal de la covarianza e innovación—, más los
registros agregados de las \num{50} corridas de Monte Carlo, del barrido de
degradación y de la malla de sensibilidad. Se incluye además un archivo JSON con
la parametrización completa y la versión de cada biblioteca del entorno de
ejecución.

\textbf{Relación con los objetivos.} Sustenta el criterio de validez interna y
trazabilidad de la sección~\ref{sec:calidad}, y provee la evidencia del proceso
exigida por el diseño metodológico.

\textbf{Estado.} Operativo; los archivos generados se adjuntan como anexo
digital.
